\chapter{Various}

\section{Intervals}
	\kactlimport{IntervalMap.h}
	\kactlimport{IntervalCover.h}
	% \kactlimport{ConstantIntervals.h}

\section{Misc. algorithms}
	\kactlimport{TernarySearch.h}
	% \kactlimport{LIS.h}
	% \kactlimport{FastKnapsack.h}

\section{Dynamic programming}
	% \kactlimport{KnuthDP.h}
	\kactlimport{DivideAndConquerDP.h}
	\subsection{Slope trick}
		To keep a convex piecewise-linear $f(x)$ with integer slopes, store the points where the slope changes by 1: left ones in a max-heap $L$, right ones in a min-heap $R$, and the minimum value $m$.
		Add $|x - a|$: push $a$ to both heaps, fix the order of the tops, and add the difference to $m$ if they cross.
		Add $\max(0, x - a)$: push $a$ to $L$, move the top of $L$ to $R$.
		Add $\max(0, a - x)$: push $a$ to $R$, move the top of $R$ to $L$.
		Prefix minimum: clear $R$. Shift by $c$: lazy offsets for $L$ and $R$.
	\subsection{Sum over subsets (SOS DP)}
		Bits are numbered from 0 on the right. Let $S(x,i)$ be the set of submasks $y \subseteq x$ that equal $x$ on all bits $\ge i$ (only the lowest $i$ bits differ),
		e.g.\ $S(1011010_2, 4)$ is the pattern $101?0?0$ with $?$ free: $1010000,1010010,1011000,1011010$.
		Then $S(x,0) = \{x\}$, $S(x,n)$ is all submasks of $x$, and
		$S(x,i+1) = S(x,i)$ if bit $i$ of $x$ is 0, else $S(x,i) \cup S(x \oplus 2^i,i)$ (disjoint: bit $i$ of $y$ is 1 or 0).
		So $dp[x][i] = \sum_{y \in S(x,i)} a[y]$ has $dp[x][0] = a[x]$ and $dp[x][i+1] = dp[x][i] + [x_i = 1]\, dp[x \oplus 2^i][i]$.
		Layer $i$ only reads entries with bit $i$ clear, which it does not change, so one array is enough: $O(n 2^n)$ time, $O(2^n)$ memory.
		\\ \verb@rep(i,0,n) rep(x,0,1<<n)@ \\ \verb@  if (x >> i & 1) f[x] += f[x ^ 1 << i];@ \\
		Supersets: test \verb@~x >> i & 1@ and add \verb@f[x ^ 1 << i]@ (bit $i$ is clear in $x$, so this is $x$ with bit $i$ set). Inverse (Mobius): same loops with \verb@-=@.
		Any other operation works too (max, min, \verb@|@), since the unions are disjoint.
		Pairs with $y \mathbin{\&} z = 0$: $z$ is a submask of $\sim y$.
\section{Debugging tricks}
	\begin{itemize}
		\item \verb@signal(SIGSEGV, [](int) { _Exit(0); });@ converts segfaults into Wrong Answers.
			Similarly one can catch SIGABRT (assertion failures) and SIGFPE (zero divisions).
			\verb@_GLIBCXX_DEBUG@ failures generate SIGABRT (or SIGSEGV on gcc 5.4.0 apparently).
		\item \verb@feenableexcept(29);@ kills the program on NaNs (\texttt 1), 0-divs (\texttt 4), infinities (\texttt 8) and denormals (\texttt{16}).
	\end{itemize}

\section{Optimization tricks}
	\verb@__builtin_ia32_ldmxcsr(40896);@ disables denormals (which make floats 20x slower near their minimum value).
	\subsection{Bit hacks}
		\begin{itemize}
			\item \verb@x & -x@ is the least bit in \texttt{x}.
			\item \verb@for (int x = m; x; ) { --x &= m; ... }@ loops over all subset masks of \texttt{m} (except \texttt{m} itself).
			\item \verb@c = x&-x, r = x+c; (((r^x) >> 2)/c) | r@ is the next number after \texttt{x} with the same number of bits set.
		\end{itemize}
	\subsection{Pragmas}
		\begin{itemize}
			\item \lstinline{#pragma GCC optimize ("Ofast")} will make GCC auto-vectorize loops and optimizes floating points better.
			\item \lstinline{#pragma GCC target ("avx2")} can double performance of vectorized code, but causes crashes on old machines.
			\item \lstinline{#pragma GCC optimize ("trapv")} kills the program on integer overflows (but is really slow).
		\end{itemize}
	% \kactlimport{FastMod.h}
	% \kactlimport{FastInput.h}
	% \kactlimport{BumpAllocator.h}
	% \kactlimport{SmallPtr.h}
	% \kactlimport{BumpAllocatorSTL.h}
	% \kactlimport{Unrolling.h}
	% \kactlimport{SIMD.h}
